%======================================================================
\section{Counting approximately commuting copies of \texorpdfstring{$S_3$}{S3}}\label{sec:criterion}
%======================================================================

Every lower bound in this paper comes from one principle, proved in~\cite{DouglasMghirbiA}. If finitely many relations force many displaced
copies of $S_3$ to commute cheaply, every model of a fixed chunk must be large, because approximately commuting copies of $S_3$ carry
entropy. We fix the conventions used in both halves of the paper, state the principle in the form used here, with $a,b$ in place of the
letters $u,v$ of~\cite{DouglasMghirbiA}, and note the ceiling $\exp(O(r))$ that it can never pass. The scheme, many cheaply transported
copies of a finite group, relations whose cost is controlled by area, and a dimension bound, follows Slofstra~\cite[Sections~2--3]{Slofstra2018}.
The criterion needs only pairwise approximate commutation of displaced copies of $S_3$.

\paragraph{Conventions.} Functions are composed from right to left, groups act on the left, $[x,y]=xyx^{-1}y^{-1}$, and $\log$ is to base
two. For a finite set $S$ and a finite set $R$ of words in $S^{\pm1}$, $\Area_R(w)$ is the least number of conjugates of words of $R^{\pm1}$
whose product is $w$ in the free group on $S$, and $\infty$ if there is none. A \emph{chunk} of a group is a finite subset $E\ni1$ with the
partial multiplication $xy=z$ whenever $x,y,z\in E$ and $xy=z$ in the group. Its \emph{sofic profile} $\sigma_E(r)$ is the least $n$ for which
some $f:E\to\Sym_n$ with $f(1)=\mathrm{id}$ multiplies and separates correctly up to $1/r$: $d(f(x)f(y),f(xy))\le1/r$ whenever $xy$ is defined in
$E$, and $d(f(x),f(y))\ge1-1/r$ for distinct $x,y\in E$~\cite{Cornulier2013}. Here $d$ is the normalized Hamming distance, the fraction of points
on which two permutations differ. The normalized Hilbert--Schmidt norm of a $D\times D$ matrix is $\|X\|_2=(\Tr(X^*X)/D)^{1/2}$.

The criterion is Theorem~8.1 of~\cite{DouglasMghirbiA}, and Appendix~\ref{app:imported} reproduces its proof.

\begin{theorem}[Counting criterion]\label{thm:criterion}
Let $G$ be a group, $S\subset G$ a finite set containing $a,b$, and $R$ a finite set of words in $S^{\pm1}$, each representing $1$ in $G$, with
$a^2,b^2,(ab)^3\in R$ and $ab\ne1$ in $G$. Let $l$ be the largest length of a word in $R$, and let $E$ be the chunk consisting of $1$, $S^{\pm1}$,
the prefixes of the words of $R$ and of $ab$, and the inverses of these elements. Let $w_1,\dots,w_N$ be words in $S^{\pm1}$ and $\Lambda\ge1$ such
that
\[
\Area_R\bigl([w_ixw_i^{-1},w_jyw_j^{-1}]\bigr)\le\Lambda\qquad(i\ne j,\ x,y\in\{a,b\}).
\]
Then $\log\sigma_E(r)\ge N/16$ for every $r\ge10^8N\Lambda(2l-1)$. Moreover, for $0<\Delta\le1$ and
$K_\Delta=\sqrt{80000\,(1+\log(1/\Delta))}/\Delta$, every assignment $\phi$ of unitaries in $U(D)$ to $S$ under which every word of $R$ lies
within $e$ of $I$ in the normalized Hilbert--Schmidt norm and $\|\phi(ab)-I\|_2\ge\Delta$ satisfies $\log D\ge N\Delta^2/12$ whenever
$K_\Delta\sqrt N\,\Lambda e\le1$.
\end{theorem}

The proof in~\cite{DouglasMghirbiA} has four steps, and only the last is specific to $S_3$. Conjugating the images of $a$ and $b$ by the image
of $w_i$ gives $N$ copies of $S_3$, one for each word. Each copy satisfies the relations of $S_3$ up to the defect of a relator, because those
relations are conjugates of $a^2$, $b^2$ and $(ab)^3$, so it can be rounded to an exact representation of $S_3$. Its three-cycle is a conjugate
of the image of $ab$, so it is displaced from the identity as much as that image is. Two copies commute up to $\Lambda$ times the relator defect,
because each of their commutators is a product of at most $\Lambda$ conjugates of relators. Finally, exact copies of $S_3$ that nearly commute
and have displaced three-cycles cannot share a small space: they force a fixed amount of entropy per copy, at least $\Delta^2/12$ bits for
unitaries and $1/16$ bit for permutations~\cite[Section~4]{DouglasMghirbiA}. The chunk $E$ does not depend on $N$ or on the words $w_i$; only
the areas matter.

The theorem has a ceiling. Its hypotheses force $N\le r/10^8$ and $N\le(K_\Delta e)^{-2}$, so it never gives $\log\sigma_E(r)$ beyond order $r$,
nor $\log D$ beyond order $e^{-2}$. Below the ceiling, everything depends on how cheaply the group carries copies of $S_3$ to many places.
Transport by the powers of one element, with area of order $k^p$ at distance $k$, gives exponent $1/(p+1)$; this is Corollary~8.2
of~\cite{DouglasMghirbiA}.

\begin{corollary}[Transport by the powers of one element]\label{cor:cyclic}
Let $G$, $S$, $R$, $l$ and $E$ be as in Theorem~\ref{thm:criterion}, let $t\in S$, and suppose that
$\Area_R([x,t^kyt^{-k}])\le Ck^p$ for $x,y\in\{a,b\}$ and $k\ge1$, with $C\ge1$ and $p\ge0$. Then
\[
\log\sigma_E(r)\ \ge\ \frac1{16}\Big\lfloor\Big(\frac r{10^8C(2l-1)}\Big)^{1/(p+1)}\Big\rfloor\qquad(r>1),
\]
and every unitary assignment as in Theorem~\ref{thm:criterion} has $\log D\ge\frac{\Delta^2}{12}\bigl((CK_\Delta e)^{-2/(2p+1)}-1\bigr)$.
\end{corollary}

In Houghton's group $H_3$ the copies of $S_3$ along a ray have commutators of at least quadratic area~\cite[Proposition~3.4]{DouglasMghirbiA},
so transport along a ray never passes the exponent $1/3$. A group that does better must reach many sites cheaply.

One more fact of Cornulier's is used twice, in Sections~\ref{sec:more} and~\ref{sec:ladder}. The profile of a chunk is bounded exactly when the
chunk embeds in a finite group, and otherwise it is at least $r$~\cite[Fact~3.5]{Cornulier2013}, indeed at least
$2r$~\cite[Lemma~2.1]{DouglasMghirbiA}. A chunk with unbounded profile therefore certifies that its group is not locally embeddable into finite
groups, not LEF~\cite[Fact~3.8]{Cornulier2013}.
